\documentclass[11pt]{article}

\usepackage[margin=1.08in]{geometry}
\usepackage{amsmath,amssymb,amsthm,mathtools}
\usepackage[T1]{fontenc}
\usepackage{lmodern}
\usepackage{microtype}
\usepackage{xcolor}
\usepackage{hyperref}
\hypersetup{
  colorlinks=true,
  linkcolor=blue!55!black,
  citecolor=blue!55!black,
  urlcolor=blue!55!black,
  pdftitle={The generic-degree spectrum of Keller counterexamples in dimension three}
}

\newtheorem{theorem}{Theorem}
\newtheorem{lemma}[theorem]{Lemma}
\newtheorem{corollary}[theorem]{Corollary}
\theoremstyle{remark}
\newtheorem{remark}[theorem]{Remark}
\newcommand{\A}{\mathbb A}
\newcommand{\C}{\mathbb C}
\newcommand{\Frac}{\operatorname{Frac}}
\newcommand{\Sing}{\operatorname{Sing}}

\title{\bfseries The generic-degree spectrum of Keller
counterexamples in dimension three\\[3pt]
\large An audited, non-novel record}
\author{Alec Kriebel, with substantial AI assistance}
\date{Program record: 24 July 2026, 22:42:23 UTC\\
First public release: 24 July 2026, 23:04:08 UTC}

\begin{document}
\maketitle

\begin{center}
\fcolorbox{black!45}{yellow!8}{%
\begin{minipage}{0.90\textwidth}
\small
\textbf{Status, priority, and review notice.}
The degree-two exclusion is classical: Campbell proved the complex Galois
case in 1973, followed by characteristic-zero algebraic proofs of Razar and
Wright.  The minimum generic degree three and the every-degree weighted-lift
family were public before this record.  No novelty is claimed.  This note is
not peer reviewed.  Exact computations check encoded identities only; they
are evidence, not peer review or verification of the universal geometric
arguments.
\end{minipage}}
\end{center}

\begin{abstract}
We record the exact generic-degree spectrum of Keller counterexamples
$\A^3_\C\to\A^3_\C$.  Generic degrees one and two cannot occur, by the
classical theorem that a Keller map with Galois function-field extension is
an automorphism.  We also give a short independent proof of the quadratic
case: a square-root generator would force a hypersurface into the
codimension-two singular locus of a reduced hypersurface.  Existing explicit
weighted-lift maps realize every generic degree at least three.  Hence the
spectrum is exactly $\{3,4,5,\ldots\}$.  Every cubic counterexample has
geometric monodromy $S_3$.  This note is an attributed research-program
checkpoint, not a new result.
\end{abstract}

\section{Statement and scope}

Let $k$ be a field of characteristic zero and let
\[
 F=(F_1,\ldots,F_n):\A^n_k\longrightarrow\A^n_k
\]
be polynomial with $\det JF\in k^\times$.  Put
\[
 A=k[F_1,\ldots,F_n]\subset B=k[x_1,\ldots,x_n],\qquad
 K=\Frac A\subset L=\Frac B.
\]
The \emph{generic degree} is $[L:K]$.

\begin{theorem}[Quadratic exclusion]\label{thm:quadratic}
If $[L:K]\leq2$, then $F$ is a polynomial automorphism.  Thus a
characteristic-zero Keller counterexample has neither generic degree one nor
generic degree two.
\end{theorem}

\begin{theorem}[Exact complex spectrum]\label{thm:spectrum}
The generic degrees of Keller counterexamples
$\A^3_\C\to\A^3_\C$ are exactly
\[
 \{3,4,5,\ldots\}.
\]
Every generic-degree-three counterexample has geometric monodromy $S_3$.
\end{theorem}

These statements concern function-field degree, not total polynomial degree
$\max_i\deg F_i$.  The latter is a separate problem.  Vistoli's theorem in
dimension three and total degree three gives the presently certified
total-degree floor $4$.

\section{The classical Galois proof}

The Keller condition makes the $F_i$ algebraically independent.  A
polynomial relation among them, after differentiation and multiplication by
the inverse of $JF$, would have zero gradient; in characteristic zero it is
constant.  Thus $F$ is dominant and $L/K$ is finite and separable.

Campbell--Razar--Wright prove the following dimension-free statement.

\begin{theorem}[Galois case {\cite{Campbell,Razar,Wright}}]\label{thm:galois}
If $F:\A^n_k\to\A^n_k$ is Keller in characteristic zero and $L/K$ is
Galois, then $A=B$.
\end{theorem}

The theorem requires neither properness nor surjectivity.  A degree-one
extension is trivially Galois, and every separable quadratic extension is
Galois.  This proves Theorem~\ref{thm:quadratic}.

\begin{remark}
The hypothesis is that the original extension $L/K$ is Galois.  The existence
of a Galois closure is not enough: every finite separable extension has one.
\end{remark}

\section{An independent quadratic argument}

Over $\C$ the quadratic case has a short proof that exposes the codimension
obstruction directly.

\begin{proof}[Second proof of Theorem~\ref{thm:quadratic} over $\C$]
Suppose $[L:K]=2$.  Since
$A\simeq\C[y_1,\ldots,y_n]$ is a unique factorization domain, the nontrivial
square class can be represented by a nonconstant squarefree $h\in A$:
\[
 L=K(\sqrt h).
\]
Set $g=\sqrt h\in L$.  It satisfies
\[
 g^2=h(F)\in B.
\]
Thus $g$ is integral over the integrally closed ring $B$, and consequently
$g\in B$.

The polynomial $g$ is nonconstant, so it has a zero.  At every $p\in V(g)$,
\[
 0=d(g^2)_p=d(h\circ F)_p=(dh)_{F(p)}\circ dF_p.
\]
Since $dF_p$ is invertible, $dh$ vanishes at $F(p)$; also $h(F(p))=0$.
Therefore
\begin{equation}
 F(V(g))\subseteq \Sing V(h).                         \label{eq:sing-image}
\end{equation}
The map $F$ is etale and hence quasi-finite.  Each irreducible hypersurface
component of $V(g)$ has image closure of dimension $n-1$.  But squarefreeness
of $h$ in characteristic zero makes $V(h)$ reduced, so it has no
codimension-one singular component:
\begin{equation}
 \dim\Sing V(h)\leq n-2.                              \label{eq:sing-codim}
\end{equation}
Equations~\eqref{eq:sing-image} and~\eqref{eq:sing-codim} contradict each
other.
\end{proof}

This proof does not extend formally to a prime degree $p>2$.  Such an
extension need not be Galois and need not be generated by a $p$th root that
integrality places in the source polynomial ring.

\section{Every generic degree at least three}

We record an explicit existing weighted-lift family.  Fix $d\geq3$ and put
\[
 u=1+xy,\qquad
 \gamma=1-\frac d{d-1}xy+x^2z.
\]
Define $F_d=(A_d,B_d,C_d)$ by
\begin{align*}
 A_d&=\frac{(d-2)u+u^2-(d-1)u^d\gamma^{d-2}}
              {(d-2)x^2},\\
 B_d&=\frac{(d-2)+2u-du^{d-1}\gamma^{d-2}}
              {(d-2)x},&
 C_d&=x\gamma.                                       \tag{3}
\end{align*}
For a direct divisibility check, put $v=xy$ and $\tau=x^2z$.  Both numerators
are polynomials in $v,\tau$; the numerator of $B_d$ has zero constant term,
whereas that of $A_d$ has zero constant and zero $v$-linear term.  Every
remaining monomial is divisible by $x$ and $x^2$, respectively.  Hence (3)
defines a polynomial map over $\mathbb Q$.

Set $w=u\gamma$, $P=B_dC_d$, and $Q=A_dC_d^2$.  Elimination gives
\[
 w^d-w^2+(d-2)Pw-(d-2)Q=0.                           \tag{4}
\]
After $U=(d-2)P$ and $V=-(d-2)Q$, the inverse equation becomes
\[
 T^d-T^2+UT+V.                                       \tag{5}
\]
It is irreducible over $\C(U,V)$: as an element of
$\C[U,V,T]$ it is primitive and linear in $V$, and any factor independent of
$V$ would have to divide its coefficient $1$.  Conversely, away from a
proper denominator locus, a root $w$ recovers the source through
\[
 \gamma=P-\frac{2w-dw^{d-1}}{d-2},\quad
 x=\frac{C_d}{\gamma},\quad u=\frac w\gamma,\quad
 y=\frac{u-1}{x},\quad
 z=\frac{\gamma-1+\frac d{d-1}(u-1)}{x^2}.           \tag{6}
\]
Thus there is no branch introduced by denominator clearing, and the generic
degree is exactly $d$.

For the Jacobian calculation, write
\[
 p_d(w)=\frac{2w-dw^{d-1}}{d-2},\qquad
 q_d(w)=\frac{w^2-(d-1)w^d}{d-2}.
\]
The identity $q_d'(w)=wp_d'(w)$ collapses the $2\times2$ Jacobian after the
target change
\[
 (A_d,B_d,C_d)\mapsto(B_dC_d,A_dC_d^2,C_d).
\]
Comparison with the input changes
$(x,y,z)\mapsto(x,xy,x^2z)\mapsto(x,w,\gamma)$ gives
\[
 \det JF_d=1.                                        \tag{7}
\]

It remains to certify noninjectivity.  Let
\[
 s_d=\frac{4-2^d}{d-2}.
\]
Over the target $(A,B,C)=(s_d,s_d,1)$, equation (4) has the distinct roots
$1$ and $2$.  The recovery denominators are
\[
 \gamma_1=\frac{d+2-2^d}{d-2}\neq0,\qquad
 \gamma_2=2^{d-1}\neq0.
\]
Equations (6) therefore give two distinct rational source points with the
same image.  So $F_d$ is a counterexample of generic degree $d$.
Together with Theorem~\ref{thm:quadratic}, this proves
Theorem~\ref{thm:spectrum}.

\section{Monodromy}

A transitive subgroup of $S_3$ is either $C_3$ or $S_3$.  If a separable
cubic extension had geometric monodromy $C_3$, it would itself be Galois.
Theorem~\ref{thm:galois} would then make the Keller map an automorphism.
Hence every cubic counterexample has geometric monodromy $S_3$.

More generally, a counterexample cannot have regular monodromy.  In a regular
transitive action the point stabilizer is trivial, so the original extension
is already Galois.

\section{Verification and disclosure}

The universal exclusion was checked against methodologically independent
proofs: Campbell's complex-analytic covering/Hartogs argument, Wright's
algebraic characteristic-zero argument, the singular-locus proof reproduced
above, and Trushin's 2026 branch-divisor theorem.  Companion SymPy and PARI/GP
scripts regression-test (3), (7), and the explicit collision for
$3\leq d\leq8$.  Finite computation is not a proof for all $d$.

This record was produced with substantial AI assistance in theorem search,
proof reconstruction, adversarial checking, exact computation, writing, and
typesetting.  It has not been peer reviewed and should not be represented as
an independent discovery.

\begin{thebibliography}{9}

\bibitem{Campbell}
L.~A. Campbell,
\emph{A condition for a polynomial map to be invertible},
Math. Ann. \textbf{205} (1973), 243--248.
\href{https://doi.org/10.1007/BF01349234}{doi:10.1007/BF01349234}.\par\smallskip

\bibitem{Razar}
M. Razar,
\emph{Polynomial maps with constant Jacobian},
Israel J. Math. \textbf{32} (1979), 97--106.
\href{https://doi.org/10.1007/BF02764906}{doi:10.1007/BF02764906}.\par\smallskip

\bibitem{Wright}
D. Wright,
\emph{On the Jacobian conjecture},
Illinois J. Math. \textbf{25} (1981), 423--440.
\href{https://doi.org/10.1215/ijm/1256047158}{doi:10.1215/ijm/1256047158}.\par\smallskip

\bibitem{Trushin}
A. Trushin,
\emph{Contracted divisors and Degree-Two Maps},
\href{https://arxiv.org/abs/2605.26390}{arXiv:2605.26390} (2026).\newline

\bibitem{Vistoli}
A. Vistoli,
\emph{The Jacobian conjecture in dimension 3 and degree 3},
J. Pure Appl. Algebra \textbf{142} (1999), 79--89.
\href{https://doi.org/10.1016/S0022-4049(98)00040-1}
{doi:10.1016/S0022-4049(98)00040-1}.\par\smallskip

\bibitem{Gallagher}
A. Gallagher,
\emph{An infinite family of counterexamples to the Jacobian Conjecture in
dimension three: every generic fiber degree $n\geq3$ occurs},
public construction (20 July 2026);
\url{https://jacobianfun.org/jacobian-explained}.

\end{thebibliography}

\end{document}
